\PassOptionsToPackage{sols}{handout}
\documentclass[main]{subfiles}
\setcounterrefbeforedot{chapter}{m-ws1-5}
\setcounterrefafterdot{section}{m-ws1-5}
\addtocounter{section}{-1}
\usepackage{solutions}

\begin{document}

\section{Logarithmic Differentiation} \label{ws1-5}

% \begin{objectives}

% Learning Objectives:

% \begin{itemize}[topsep=0pt,partopsep=0pt,parsep=8pt,itemsep=0pt]

% \item You should be able to use logarithmic differentiation to find the
%   derivatives of functions of the form $f(x)^{g(x)}$.

% \item You should recognize other situations in which logarithmic
%   differentiation may be simpler than other methods of differentiation
%   (for instance, if you have a product or quotient involving lots of
%   terms).

% \end{itemize}
% \end{objectives}

\begin{enumerate}
\item  On your homework you looked at $y= x^x$.  This is
  neither an exponential function (it's not of the form $b^x$ with $b$ a
  constant) nor a power function (it's not of the form $x^n$ with $n$ a
  constant).  You showed that $y'$ is
  neither $\ln x \cdot x^x$ nor $x \cdot x^{x-1}$.  In this problem, we'll
  look at two ways of finding $y'$. 

  \begin{enumerate}
  \item

    \note{Here, I'll probably talk a bit about how a useful problem
      solving strategy is to ``transform'' a problem you don't know how to
      do into a problem you do know how to do.
      
      Matt note - I'll do this with them and specifically label Robin's note above as a helpful "thinking strategy"}

    \begin{problem}[\stretch{2}]
      \textbf{Method 1.}  Rewrite $y = x^x$ as
      $y=e^{g(x)}$. Simplify and use the Chain Rule to find
      $y'$.
    \end{problem}

    \begin{solution}
      Since $e^x$ and $\ln x$ are inverse functions, we can write
      \begin{align*}
        x^x &= e^{\ln (x^x)} \\
        \intertext{Using log rules, we can simplify this as} &= e^{x \ln
          x}
      \end{align*}
      Now, we can use the Chain Rule to differentiate this:
      \begin{align*}
        \frac{d}{dx} \left( x^x \right) &= \frac{d}{dx}
        \left(\fcolorbox{red}{white}{$\displaystyle e^ {
              \fcolorbox{blue}{white}{$\displaystyle x \ln x$}}
            $}\right) \\
        &= e^ { x \ln x} \cdot \frac{d}{dx}
        \left(\fcolorbox{green}{white}{$\displaystyle x$}\cdot
          \fcolorbox{green}{white}{$\displaystyle  \ln x$}\right) \\
        &= e^ { x \ln x} \cdot \left(x\cdot
          \frac{1}{x}+ \ln x\right) \\
        &= e^{x \ln x} (1 + \ln x) \\
        \intertext{Finally, remember that $e^{x \ln x} = x^x$.}
      \end{align*} 
    \end{solution}

  % \item

  %   \note{I want to make a couple of points here.  The first is that we
  %     really should only use logarithmic differentiation when the function
  %     is positive; the second is that it's not an issue in this case,
  %     since the domain of this function is essentially positive numbers,
  %     and then the function is positive.  (I won't go in-depth on the
  %     domain, but I'll ask them something like whether the function makes
  %     sense with $x = - \dfrac{1}{2}$, and they should see that it
  %     doesn't.)} 

    % \begin{problem}[\stretch{0.5}]
    %   This only makes sense when $x^x$ is positive.  What part goes
    %   wrong when $x^x \leq 0$?  Should we worry about this?
    % \end{problem}

    % \begin{solution}
    %   In the very first step, we rewrote $x^x = e^{\ln (x^x)}$, and
    %   $\ln(x^x)$ is only defined when $x^x$ is positive (because we can
    %   only take the natural log of a positive number).

    %   However, this is not an issue.  Let's think about possible inputs
    %   to $f(x) = x^x$.  Certainly, it makes sense to take $f$ of any
    %   positive number, but it doesn't make sense to take $f$ of most
    %   negative numbers; for example $f \left( - \frac{1}{2} \right) =
    %   \left( - \frac{1}{2} \right)^{-1/2} =
    %   \frac{1}{\sqrt{-\frac{1}{2}}}$, which is undefined.  So it makes sense to only consider the domain of $f(x) = x^x$ to be $(0, \infty)$, in which case there's no issue.
    % \end{solution}

  \item \note{In Ma, we did (in rare instances) differentiate both sides
      of an equation, but we never spent much time talking about it. We
      don't need to belabor the point here, but it's worth pointing out to
      students that taking the derivative of both sides of an equation is
      different from what we did in Method 1. This is part of the bigger
      theme that simplifying an expression and performing an operation on
      both sides of an equation are two different things.}

    \textbf{Method 2 (logarithmic differentiation).}  
    \begin{enumerate}
    \item 
    \begin{problem}[\vfill]
      Take the natural log of both sides of $y=x^x$ and simplify the right side using log rules.
    \end{problem}
    
    \solonly{\[\ln y = \ln (x^x) = x \ln x\]}

    \item
    \begin{problem}[\vfill]
      Differentiate both sides with respect to $x$.
    \end{problem}

    \solonly{\[\frac{y'}{y} = 1 \cdot
      \ln x + x \cdot \frac{1}{x} = \ln x + 1\]}

    \item
    \begin{problem}[\vfill]
      Solve for $y'$. Replace $y$ with $x^x$ to get an explicit formula for $y'$.
    \end{problem}

    \solonly{\[f'(x) = f(x) (\ln x + 1) = x^x (\ln x + 1)\]}
    \end{enumerate}

  \item
    \begin{problem}
      Do the two methods give you the same answer?  Which do you like
      better? 
    \end{problem}
    \pspace{\vfill}

    \begin{solution}
      The two methods definitely give the same answer; which you prefer is
      up to you!
    \end{solution}

  \end{enumerate}
\end{enumerate}

\begin{framed}
    \textbf{Logarithmic Differentiation}
    \begin{enumerate}
        \item Take the natural log of both sides of an equation $y=f(x)$ and use log rules to simplify.
        \item Differentiate both sides of the equation with respect to $x$.
        \item Solve for $y'$ and substitute $f(x)$ for $y$. 
    \end{enumerate}
\end{framed}

\newpage

\textbf{Key Idea}: We can use logarithmic differentiation to differentiate $y=f(x)^{g(x)}$ (variable base and exponent.)

% \note{Again, encourage the students to quiz themselves here.  You may want
%   to talk a little about how they can reconstruct these if they've
%   forgotten them.}

\begin{enumerate}[resume]
\item  In each part, find the derivative.

  \begin{enumerate}

  \item
    \begin{problem}[\stretch{1}]
      $y = (4x^2 + 5)^{e^x-1}$.      
    \end{problem}

    \begin{solution}
      We need logarithmic differentiation because both the base and
      exponent depend on $x$.  We start with the given equation: 
      \begin{align*}
        y &= (4x^2 + 5)^{e^x-1} \\
        \intertext{Take the natural log of both sides and simplify the
          right side:}
        \ln y &= \ln \left[ (4x^2 + 5)^{e^x - 1} \right] \\
        &= (e^x - 1) \ln (4x^2 + 5) \text{ by {L6}} \\
        \intertext{Differentiate both sides:} %
        \frac{y'}{y} &= \frac{d}{dx} \left(\fcolorbox{green}{white}{$\displaystyle  \left(e^ { x} -1\right)$}\cdot \fcolorbox{green}{white}{$\displaystyle  \ln\left(4 x^ { 2} + 5\right)$}\right) \\
        &= \left(e^ { x} -1\right)\cdot \frac{d}{dx} \left(\fcolorbox{red}{white}{$\displaystyle \ln\left(\fcolorbox{blue}{white}{$\displaystyle 4 x^ { 2} + 5$}\right)$}\right)+ \frac{d}{dx}\left(e^ { x} -1\right)\cdot \ln\left(4 x^ { 2} + 5\right) \\
        &= \left(e^ { x} -1\right)\cdot \frac{1}{4 x^ { 2} + 5}\cdot 8 x +
        e^ { x} \cdot \ln\left(4 x^ { 2} + 5\right) \\
        \shortintertext{Solve for $y'$:} %
        y' &= y \left[ \left(e^ { x} -1\right)\cdot \frac{1}{4 x^ {
              2} + 5}\cdot 8 x + e^ { x} \cdot \ln\left(4 x^ { 2} +
            5\right) \right] \\
        &= \boxed{(4x^2 + 5)^{e^x-1} \left[ \left(e^ { x} -1\right)\cdot
            \frac{1}{4 x^ { 2} + 5}\cdot 8 x + e^ { x} \cdot \ln\left(4 x^
              { 2} + 5\right) \right]}
      \end{align*}
    \end{solution}

  \item

    \note{Students will probably initially try to just take ln of both
      sides.} 

    \begin{problem}[\stretch{1}]
      $f(t) = 3^t + t^3 + t^{3t}$.      
    \end{problem}

    \begin{solution}
      We know the derivatives of $3^t$ and $t^3$, but we need logarithmic
      differentiation to find the derivative of $t^{3t}$ since the base
      and exponent both depend on $t$.  Let's give this function a new name,
      like $g(t)$:
      \begin{align*}
        g(t) &= t^{3t} \\
        \intertext{Take the natural log of both sides and simplify the
          right side:}
        \ln g(t) &= \ln \left( t^{3t} \right) \\
        &= 3t \ln t \\
        \intertext{Differentiate both sides:} %
        \frac{g'(t)}{g(t)} &= 3t \cdot \frac{1}{t} + 3 \ln t \\
        &= 3 + 3 \ln t
        \intertext{Solve for $g'(t)$:}
        g'(t) &= t^{3t} (3 + 3 \ln t) 
      \end{align*}
      So, $f'(t) = \boxed{3^t \ln 3 + 3t^2 + t^{3t} (3 + 3 \ln t)}$.
    \end{solution}

  \item
    \begin{problem}[\nstretch{1}]
      $f(x) = (\ln x)^x - \sqrt{x^2 + 1}$
    \end{problem}

    \begin{solution}
      We need logarithmic differentiation to find the derivative of $(\ln
      x)^x$, but finding the derivative of $\sqrt{x^2 + 1}$ just involves
      the Chain Rule.  So, let's give $(\ln x)^x$ a new name, like $g(x)$:
      \begin{align*}
        g(x) &= (\ln x)^x \\
        \intertext{Take the natural log of both sides and simplify the
          right side:}
        \ln g(x) &= \ln \left[ (\ln x)^x \right] \\
        &= x \ln (\ln x) \\
        \intertext{Differentiate both sides:} %
        \frac{g'(x)}{g(x)} &=
        \frac{d}{dx} \left(\fcolorbox{green}{white}{$\displaystyle
            x$}\cdot \fcolorbox{green}{white}{$\displaystyle  \ln\left(\ln
              x\right)$}\right) \\
        &= x\cdot \frac{d}{dx} \left(\fcolorbox{red}{white}{$\displaystyle
            \ln\left(\fcolorbox{blue}{white}{$\displaystyle \ln
                x$}\right)$}\right)+ \ln\left(\ln x\right) \\ 
        &= x\cdot \left(\frac{1}{\ln x}\cdot \frac{1}{x}\right)+
        \ln\left(\ln x\right) \\
        &= \frac{1}{\ln x} + \ln(\ln x)\\
        \intertext{Solve for $g'(x)$:} g'(x) &= (\ln x)^x \left(
          \frac{1}{\ln x} +\ln(\ln x) \right)
      \end{align*}
      Therefore,
      \begin{align*}
        f'(x) = {} & (\ln x)^x \left(
          \frac{1}{\ln x} +\ln(\ln x) \right) - \frac{d}{dx}
        \left(\fcolorbox{red}{white}{$\displaystyle
            \fcolorbox{blue}{white}{$\displaystyle  \left(x^ { 2} +
                1\right)$}^ { \frac{1}{2}} $}\right) \\
        = {} & (\ln x)^x \left(
          \frac{1}{\ln x} +\ln(\ln x) \right) - \frac{1}{2}\cdot \left(x^
          { 2} + 1\right)^ { -\frac{1}{2}} \cdot \left(2 x\right) 
      \end{align*}
    \end{solution}

  \end{enumerate}

\end{enumerate}

\textbf{Key Idea}: Logarithmic differentiation is useful for dealing with complicated products, quotients, and powers.

\begin{enumerate}[resume]

  \item
    \note{A few common mistakes come up in simplifying $\ln y$: students
      often try to simplify $\ln(xe^x)$ as $x \ln (xe)$ (they are trying
      to ``bring the exponent out''), and they often make sign errors when
      simplifying the denominator (writing $-4 \ln (x^2 + 2) + 2 \ln (5x +
      2)$ instead of $-[4 \ln (x^2 + 2) + 2 \ln (5x + 2)]$.

      Note that this is also the first example using the notation $y$ and
      $\frac{dy}{dx}$ vs. $f(x)$ and $f'(x)$; this tends to confuse
      students when we do implicit differentiation, so you may want to
      particularly highlight the notation here. 
    }

    \begin{problem}[\stretch{1}]
      Differentiate $y = \dfrac{xe^x}{(x^2+2)^4(5x+2)^2}$.
    \end{problem}

    \begin{solution}
      We could do this without logarithmic differentiation, but then we'd
      need to do a lot of Product and Quotient Rule.  Logarithmic
      differentiation will make life easier!
      \begin{align*}
        y &= \frac{xe^x}{(x^2+2)^4 (5x+2)^2} \\
        \intertext{Take the natural log of both sides and simplify the
          right side:}
        \ln y &= \ln \left( \dfrac{xe^x}{(x^2+2)^4 (5x+2)^2} \right) \\
        &= \ln \left( xe^x \right) - \ln \left[ (x^2+2)^4 (5x+2)^2 \right]
        \text{ by {L5}} \\
        &= \ln x + \ln e^x - \ln [(x^2 + 2)^4] - \ln [(5x + 2)^2] \text{
          by {L3}} \\
        &= \ln x + x - 4 \ln (x^2 + 2) - 2 \ln(5x + 2) \text{ by
          {L6}} \\
        \intertext{Differentiate both sides:} \frac{1}{y} \frac{dy}{dx} &=
        \frac{1}{x} + 1- 4 \cdot
        \frac{2x}{x^2 + 2} - 2 \cdot \frac{5}{5x + 2} \\
        \intertext{Solve for $\dfrac{dy}{dx}$:} \frac{dy}{dx} &=
        \boxed{\frac{xe^x}{(x^2+2)^4 (5x+2)^2} \left( \frac{1}{x} + 1- 4
            \cdot \frac{2x}{x^2 + 2} - 2 \cdot \frac{5}{5x + 2} \right)}
      \end{align*}
    \end{solution}

\item 
  \begin{problem}[\stretch{1}]
    Suppose you want to find the derivative of $y=a(x) b(x) c(x) d(x)$.
    You \emph{could} do this using the Product Rule, but it's much easier
    with logarithmic differentiation.  Find a formula for the derivative.
  \end{problem}

  \begin{solution}
    Let's give $f(x) g(x) h(x) j(x)$ a name, like $p(x)$.  (After all, in
    logarithmic differentiation, it's important that we work with
    equations.) 
    \begin{align*}
      p(x) &= f(x) g(x) h(x) j(x) \\
      \intertext{Take the natural log of both sides and simplify the
        right side:}
      \ln p(x) &= \ln [f(x) g(x) h(x) j(x)] \\
      &= \ln f(x) + \ln g(x) + \ln h(x) + \ln j(x) \\
      \intertext{Differentiate both sides with respect to $x$:}
      \frac{p'(x)}{p(x)} &= \frac{f'(x)}{f(x)} + \frac{g'(x)}{g(x)} +
      \frac{h'(x)}{h(x)} + \frac{j'(x)}{j(x)} \\
      \shortintertext{Solve for $p'(x)$:}
      p'(x) &= f(x) g(x) h(x) j(x) \left( \frac{f'(x)}{f(x)} + \frac{g'(x)}{g(x)} +
        \frac{h'(x)}{h(x)} + \frac{j'(x)}{j(x)} \right) \\[4pt]
      &= \fbox{f'(x) g(x) h(x) j(x) + f(x) g'(x) h(x) j(x) + f(x) g(x) h'(x) j(x)
      + f(x) g(x) h(x) j'(x)}  
    \end{align*}
  \end{solution}

\item 
\begin{problem}
    Differentiate $\displaystyle y=\sqrt{\frac{x-1}{x^4+1}}$.
\end{problem}
\pspace{\vfill}

\begin{solution}
    We could use the Chain Rule, but logarithmic differentiation will be more efficient. First we take the natural log of both sides and simplify the right side:
    \begin{align*}
        \ln y &= \ln\sqrt{\frac{x-1}{x^4+1}} \\[4pt]
        &= \frac{1}{2}\ln\!\left(\frac{x-1}{x^4+1}\right) \\[4pt]
        &= \frac{1}{2}\left(\ln(x-1)-\ln(x^4+1)\right)\\
        \intertext{Now we differentiate both sides:}
        \frac{y'}{y} &= \frac{1}{2}\left(\frac{1}{x-1} - \frac{4x^3}{x^4+1}\right)\\
        \intertext{Solve for $y'$:}
        y' &= \frac{1}{2}y\left(\frac{1}{x-1} - \frac{4x^3}{x^4+1}\right)\\[4pt]
        &= \fbox{$\displaystyle\frac{1}{2}\sqrt{\frac{x-1}{x^4+1}}\left(\frac{1}{x-1} - \frac{4x^3}{x^4+1}\right)$}
    \end{align*}
\end{solution}

\pspace{\newpage}


% \item 
% \begin{problem}[\stretch{1}]
%   Suppose you forgot the derivative of $y = \displaystyle \frac{f(x)}{g(x)}$. Use logarithmic differentiation to recover it.
%     \end{problem}

%     \begin{solution}
%       \ideas{Since the base and exponent both depend on $x$, we should use
%         logarithmic differentiation.}  We start with the given equation:
%       \begin{align*}
%         y &= \frac{f(x)}{g(x)} \\
%         \intertext{Take the natural log of both sides and simplify the
%           right side:}
%         \ln y &= \ln \left( \frac{f(x)}{g(x)} \right) \\
%         &= \ln (f(x)) - \ln (g(x)) \\
%         \intertext{Differentiate both sides with respect to $x$:}
%         \frac{y'}{y} &= \frac{f'(x)}{f(x)} - \frac{g'(x)}{g(x)} \\
%         \shortintertext{Solve for $y'$:}
%         y' &= y \left(\frac{f'(x)}{f(x)} - \frac{g'(x)}{g(x)}
%         \right) \\
%         &= \frac{f(x)}{g(x)} \left(\frac{f'(x)}{f(x)} - \frac{g'(x)}{g(x)}
%         \right) \\
%         &= \frac{f'(x)}{g(x)} - \frac{f(x) g'(x)}{g(x)^2} \\
%         &= \frac{f'(x) g(x) - f(x) g'(x)}{g(x)^2},
%       \end{align*}
%       which is exactly the Quotient Rule. 
%     \end{solution}

%   \item

%     \note{I just want a chance to mention that we can always use the Chain
%       Rule to deal with quotients.}

%     \begin{problem}[\nstretch{0.75}]
%       As you can see, if you forget the Quotient Rule, you can always use
%       logarithmic differentiation to figure it out again.  What's another
%       way to figure out the Quotient Rule if you've forgotten it?
%     \end{problem}

%     \begin{solution}
%       You can also use the Chain Rule together with the Product
%       Rule. \version{Mb}{(You did this in
%         {{{Problem Set 1}, {{{{\#2}}(b)}}}}.)} 
%     \end{solution}

%   \end{enumerate}


\item
\begin{problem}[\stretch{1}]
    You are given the following information about two differentiable functions
    $g(x)$ and $h(x)$.
    \[  
    g(3)=2 \qquad g'(3)=-10 \qquad h(3)=5
    \]
    If $f(x)=g(x)^{h(x)}$, do you have enough information to find $f'(3)$?
    If so, what is $f'(3)$? If not, what additional information do you need?
\end{problem}

\begin{solution}
    We don't have enough information.
    Intuitively, the rate of change of $f$ will depend on the rates of change
    of both $g$ and $h$, so we would need $h'(3)$ in addition to the given
    information to determine $f'(3)$. We can use logarithmic differentiation
    to determine a general differentiation rule. For simplicity, let's say
    $y = f(x)$, $u = g(x)$, and $v = h(x)$.
    
    First, we take the log of both sides and simplify.
    \begin{align*}
    \ln y &= \ln u^v \\
    \ln y &= v \ln u
    \end{align*}
    Now we differentiate both sides and solve for $y'$.
    \begin{align*}
    \frac{y'}{y} &= v' \cdot \ln u + v \cdot \frac{u'}{u} \\
    y' &= v' \cdot y \ln u + v y \cdot \frac{u'}{u} \\
    &= v' \cdot u^v \ln u + v u^v \cdot \frac{u'}{u} \\
    &= u^v \ln u \cdot v' + v u^{v-1} \cdot u'
    \end{align*}
    Using this rule, we can determine that 
    \[f'(3) = g(3)^{h(3)}\ln g(3) \cdot h'(3)
    + h(3) g(3)^{h(3)-1} \cdot g'(3)\]
\end{solution}



% \item
% \begin{problem}[\stretch{0.5}]
%     Suppose you forgot the derivative of $y=b^x$. Use logarithmic differentiation
%     to recover it.
% \end{problem}

% \begin{solution}
% \begin{align*}
% \ln y &= \ln b^x \\
% \ln y &= x \ln b \\
% \shortintertext{Differentiating both sides:}
% \frac{y'}{y} &= \ln b \text{ (since $\ln b$ is a constant)} \\
% y' &= y\ln b \\
% &= b^x \ln b
% \end{align*}
% \end{solution}

\item
    There's an issue with logarithmic differentiation that we've brushed under the rug: $\ln f(x)$ is not defined if $f(x)$ is zero or negative. The reason why this isn't a problem is because of the following fact.

\begin{enumerate}

\item 
\begin{problem}
    Show that if $f(x)=\ln|x|$, then $f'(x)=\frac{1}{x}$. (Hint: Rewrite $f(x)$ as a piecewise function.) 
\end{problem}
\pspace{\stretch{0.5}}

\begin{solution}
    Recall that $|x|= x$ when $x\geq0$, and $|x|=-x$ when $x\leq0$. So we can express
    $\ln|x|$ as
    \[f(x)=
    \begin{cases}
        \ln(x),\ x>0 \\
        \ln(-x),\ x<0
    \end{cases}\]
    Since the derivative of $\ln(x)$ and $\ln(-x)$ are both $\frac{1}{x}$, $f'(x)=\frac{1}{x}$ on the entire domain.
\end{solution}

\item 
\begin{problem}
Sketch the graphs of $y=\ln x$ and $y=\ln|x|$. Based on the graphs, explain how it's possible that the derivatives of these two different functions have the same formula, $y'=\frac{1}{x}$.
\end{problem}
\pspace{\stretch{0.5}}

\begin{solution}
    The graph of $y=\ln x$ is familiar to us.
    \begin{center}
        \begin{tikzpicture}
            \begin{axis}[
            xtick=\empty, ytick=\empty,
            xmin=-10, xmax=10, ymin=-3, ymax=3
            ]
                \addplot[samples=200,domain=0:10] {ln(x)};
            \end{axis}
        \end{tikzpicture}
    \end{center}
    The graph of $y=\ln|x|$ will be the same as the graph of $y=\ln x$ when $x>0$, but when $x<0$, it will be the same as graph of $y=\ln(-x)$. The graph of $y=\ln(-x)$ is a reflection of the graph of $y=\ln(x)$ over the $y$-axis.
        \begin{center}
        \begin{tikzpicture}
            \begin{axis}[
            xtick=\empty, ytick=\empty,
            xmin=-10, xmax=10, ymin=-3, ymax=3
            ]
                \addplot[samples=200,domain=-10:10] {ln(abs(x))};
            \end{axis}
        \end{tikzpicture}
    \end{center}
    The two functions' derivatives can have the same formula because the functions and their derivatives have different domains. $y=\ln x$ and $y=\ln|x|$ are the same when $x>0$. When $x<0$, $y=\ln x$ is not defined, but $y=\ln|x|$ is the same as $y=\ln(-x)$, which also has the same derivative as $\ln(x)$, $y'=\frac{1}{x}$, just on a different domain.
\end{solution}


\item
\begin{problem}
Differentiate $y=x^3$ using logarithmic differentiation. Since $x^3$ can be negative, take the absolute value of both sides before taking logs of both sides. Does this affect this rest of the process at all?
\end{problem}
\pspace{\stretch{0.5}}

\begin{solution}
    First, we take the absolute value of both sides and take the natural log of both sides.
    \begin{align*}
        |y| &= |x^3| \\
        \ln|y| &= \ln|x^3|\\
        \intertext{Now we can differentiate both sides, using the fact that $\diff*{\ln|x|}{x}=\frac{1}{x}$:}
        \frac{y'}{y} &= \frac{3x^2}{x^3} \\
        \frac{y'}{y} &= \frac{3}{x} \\
        \intertext{Now we can solve for $y'$:}
        y' &= \frac{3y}{x} \\
        &= \frac{3x^3}{x} \\
        &= \fbox{$3x^2$}
    \end{align*}
\end{solution}
We can see that taking the absolute value of both sides did not affect the rest of the process once we differentiated, and the answer was what we expected.
\end{enumerate}

\end{enumerate}

\end{document}